\documentclass[10pt]{beamer}

%%%%%%%%%%%%%%%%%%%%%%%%%
%   Options for notes   %
%%%%%%%%%%%%%%%%%%%%%%%%%
% \setbeamertemplate{note page}[plain]
% \setbeameroption{show notes}
% \setbeameroption{show notes on second screen}
% \setbeameroption{show only notes}
% \setbeameroption{show only slides with notes}

\input{preamble}

%%%%%%%%%%%%%%%%%%%%%%%%%
%     presenter info    %
%%%%%%%%%%%%%%%%%%%%%%%%%
\author[Xinda GE]{
Xinda GE\\
\small PhD year 1 \\ Guided by Prof. Stephen Dalton \& Dr. Tianming WU
}
\title[SOX2 in hESC by Degron system]{
Using Degron system to dissect the function of SOX2 in human ES cell differentiation
}
% \subtitle{
% Guided by Prof. Stephen Dalton \& Dr. Tianming WU
% }
\institute[Dalton Lab]{
Dalton Lab,\\
School of Biomedical Sciences
}
\date{\today}


\begin{document}

% \kaishu
\begin{frame}
    \titlepage
    \begin{figure}[htpb]
        \begin{center}
            
\includegraphics[width=0.3\linewidth]{pic/CUHKSZ-Logo.pdf}
        \end{center}
    \end{figure}
\end{frame}

%%%%%%%%%%%%%%%%%%% Contents Start %%%%%%%%%%%%%%%%%%%%%%
\section{Method}

\begin{frame}{Inducible Ablation of a Targeted Protein Molecule in nucleus}  
    \begin{itemize}
    \end{itemize}
        
    \begin{figure}
        \centering
        \includegraphics[width=0.9\textwidth]{20230422-SOX2/figs/Fig1.png}
        \caption{Approaches to clear protein molecules inside the nucleus of a cell}
        \label{fig:clearprotein}
    \end{figure}
\end{frame}

\begin{frame}{What is needed to create a Degron cell line?}  
    \begin{itemize}
    \end{itemize}
        
    \begin{figure}
        \centering
        \includegraphics[width=1.0\textwidth]{20230422-SOX2/figs/Fig2.png}
        \caption{AID system mechanism}
        \label{fig:mechanisms}
    \end{figure}
\end{frame}

\begin{frame}{Effectiveness Validation}  
    \begin{itemize}
    \end{itemize}
        
    \begin{figure}
        \centering
        \includegraphics[width=1.0\textwidth]{20230422-SOX2/figs/Fig3.png}
        \caption{Validations}
        \label{fig:val}
    \end{figure}
\end{frame}

\section{Assumption}

\section{Experiment Design}

\begin{frame}{Gene Editing}  
    \begin{itemize}
    \end{itemize}
        
    \begin{figure}
        \centering
        \includegraphics[width=0.9\textwidth]{20230422-SOX2/figs/Fig4.png}
        \caption{Gene Editing Process}
        \label{fig:gene editing}
    \end{figure}
\end{frame}
\section{Temporary Results}

\subsection{PCR Results}
\subsection{Western Blot Results}
\subsection{Cell Sorting Results}
\subsection{ Results}

\begin{frame}{Itemize \& Enumerate}
    \begin{itemize}
        \item this the first highlight key words
        \item this the second highlight key words
        \begin{itemize}
            \item sub-keyword 1
            \item sub-keyword 2
        \end{itemize}
    \end{itemize}
    \begin{enumerate}
        \item list 1
        \item list 2
    \end{enumerate}
\end{frame}




\begin{frame}{left:right = 3:7}
    \begin{columns}
    \column{0.3\textwidth}
        This is column one with 0.3 text width.
        All other commands can be nested in the columns env.
    \column{0.7\textwidth}
        This is column two with 0.7 text width.

        I would like to cite Figure \ref{fig:test} here, which is shown on the next page.
    \end{columns}
\end{frame}


\begin{frame}{Add Figure}
    This slide is to show how to include a figure and cite it as Figure \ref{fig:test}.
    \begin{figure}
        \centering
        \includegraphics[width=0.7\textwidth]{20230422-SOX2/figs/test.pdf}
        \caption{Test Figure}
        \label{fig:test}
    \end{figure}
\end{frame}

\begin{frame}{Add Table}
    \begin{table}
        \centering
        \caption{Test Table\footnote{Table title is usually placed on the top of table.}}
        \begin{tabular}{ccc}
            \toprule
            1st column & 2nd column & 3rd column \\
            \midrule
            line 1 & blahblah & an formula $e^{\pi i} + 1 = 0$ \\
            line 2 & xxxxxx & another formula $\begin{bmatrix}1 & 2 \\ 3 & 4\end{bmatrix}$ \\
            \bottomrule
        \end{tabular}
        \label{tab:texttab}
    \end{table}
\end{frame}




%%%%%%%%%%%%%%%%%%% Contents End %%%%%%%%%%%%%%%%%%%%%%

\section*{References}
\begin{frame}{References}
    \footnotesize
    \nocite{*}
    \bibliography{ref}
    \bibliographystyle{abbrv}
    % \tiny\bibliographystyle{alpha}
\end{frame}


\begin{frame}
    \begin{center}
        {\Huge\calligra Thanks!}
    \end{center}
\end{frame}

\end{document}